\documentclass{article}
\usepackage[margin=25mm]{geometry}
\usepackage{amsmath,amssymb}
\usepackage{xurl}
\usepackage[hidelinks]{hyperref}
\setlength{\emergencystretch}{2em}
% Shared commands for notes in docs/paper-gaps/.

\DeclareUnicodeCharacter{2080}{\ifmmode{}_{0}\else\textsubscript{0}\fi}
\DeclareUnicodeCharacter{2081}{\ifmmode{}_{1}\else\textsubscript{1}\fi}
\DeclareUnicodeCharacter{2082}{\ifmmode{}_{2}\else\textsubscript{2}\fi}
\DeclareUnicodeCharacter{2099}{\ifmmode{}_{n}\else\textsubscript{n}\fi}

\newcommand{\Lean}{\textsc{Lean}}
\ExplSyntaxOn
\NewDocumentCommand{\leanid}{m}{
  \ifmmode
    \textsf{\detokenize{#1}}
  \else
    \group_begin:
    \urlstyle{sf}
    \tl_set:Nn \l_tmpa_tl {#1}
    \tl_replace_all:Nnn \l_tmpa_tl {\_} {_}
    \exp_args:NV \path \l_tmpa_tl
    \group_end:
  \fi
}
\ExplSyntaxOff
\providecommand{\tr}{\operatorname{tr}}
\newcommand{\tnleanIssueBase}{https://github.com/LionSR/TNLean/issues/}
\newcommand{\tnleanPrBase}{https://github.com/LionSR/TNLean/pull/}
\newcommand{\tnleanDocBase}{https://sirui-lu.com/TNLean/docs/}
\newcommand{\ghissue}[1]{\href{\tnleanIssueBase#1}{GitHub issue \##1}}
\newcommand{\ghpr}[1]{\href{\tnleanPrBase#1}{PR \##1}}
\newcommand{\doclink}[1]{\href{\tnleanDocBase#1.html}{\path{#1}}}

% Dirac notation and the identity operator, as in blueprint/src/macros/common.tex.
% \providecommand keeps note-local definitions authoritative.
\usepackage{dsfont}
\providecommand{\ket}[1]{|#1\rangle}
\providecommand{\bra}[1]{\langle #1|}
\providecommand{\braket}[2]{\langle #1 | #2 \rangle}
\providecommand{\ketbra}[2]{|#1\rangle\!\langle #2|}
\providecommand{\Id}{\mathds{1}}
\providecommand{\one}{\mathds{1}}
\providecommand{\C}{\mathbb{C}}
\providecommand{\MN}[1]{M_{#1}(\C)}
\providecommand{\MD}{M_D(\C)}

% Verdict marker: \gapnote{<kind>}{<status>}. Kind is the policy
% classification (clarification, local-correction, scope-restriction,
% unfaithful, false-source, open-gap); status is open, wip (elimination
% underway), resolved, or historical. Machine-read by the site index;
% prints a verdict line.
\newcommand{\gapnote}[2]{\par\noindent\textbf{Verdict:} #1 (#2).\par}

\title{The bond dimension of the middle factors of the efficient Chebyshev projector}
\author{}
\date{2026-11-19}
\begin{document}
\maketitle
\gapnote{scope-restriction}{open}

\section{The source statement}

Proposition~5.7 (\texttt{prop:Cheby-AGSP-eff}) of Arad, Landau, Vazirani and
Vidick~\cite{Arad2016RigorousRG}, arXiv:1602.08828v2, file
\path{local_alg_arxiv_v2.tex}, lines 1311--1320, reads, with its proof at
lines 1322--1333:
\begin{quote}
Let $H$ be a Hamiltonian having a decomposition of the form described
in~\texttt{eq:H-mid-cut}, $k$ an integer, and $K = P_k(H)$ the associated
degree-$k$ Chebyshev AGSP as defined in Definition~\texttt{def:Cheby-AGSP}.
Assume that $H_M$ (but not necessarily $H_L$ or $H_R$) is specified by an MPO
with bond dimensions at most $\tilde{B}$.

Then there exists $D \leq (dk)^{O(\ell + k/\ell)}$ such that a family of $D^2$
MPO $\{A_1,\ldots,A_{D^2}\}$ of bond dimension at most $\tilde{B}^k$ each such
that there exists $B_1,\ldots,B_{D^2}$ with $K=\sum A_i\otimes B_i$ can be
computed in time $nD^2\tilde{B}^{O(k)}$. Here the $A_i$ act on qudits
$\{i_1,i_1+1,\ldots,i_2\}$ and the $B_i$ on the remaining qudits. This
computation does not require knowledge of $\eta_0,\eta_1$.

Furthermore, if $H_L$ and $H_R$ are also given as MPO with bond dimension at
most $\tilde{B}$ then the $B_i$ can be computed as well.
\end{quote}
The decomposition \texttt{eq:H-mid-cut} is the one of Theorem
\texttt{thm:Cheby-AGSP} (lines 809--834). For $i_1<i_2$ in $\{1,\ldots,n\}$
and $3\leq\ell\leq(i_2-i_1)/2$ it reads
\begin{align}
 H={}&H_L+h_{i_1-\ell}+\cdots+h_{i_1}+\cdots+h_{i_1+\ell-1}\notag\\
 &+H_M+h_{i_2-\ell}+\cdots+h_{i_2}+\cdots+h_{i_2+\ell-1}+H_R,
 \label{eq:alvv16_cmb_midcut}
\end{align}
where each $h_i$ is a two-local operator on the qudits $\{i,i+1\}$, and $H_L$,
$H_M$, $H_R$ act on
\begin{align}
 J_L=\{1,\ldots,i_1-\ell\},\qquad
 J_M=\{i_1+\ell,\ldots,i_2-\ell\},\qquad
 J_R=\{i_2+\ell,\ldots,n\}.
 \label{eq:alvv16_cmb_regions}
\end{align}
The quotation keeps the source's labelling of the sites by $1,\ldots,n$. The
running time and the computability of the $B_i$ are out of scope; only the
existence content of the proposition is considered. The two groups of terms
$h_{i_1-\ell},\ldots,h_{i_1+\ell-1}$ and $h_{i_2-\ell},\ldots,h_{i_2+\ell-1}$
are called the left and right windows. The bound $\tilde B$ concerns $H_M$
alone, not the windows.

\section{What the printed proof gives}

The proof refers to the proof of Lemma~4.2 of~\cite{Arad2013AreaLaw}. It expands $K$
through the operators $P_{u_1u_2,kj_1j_2}(Z)$, the sum of the terms of
$(H_L+\cdots+H_R)^k$ with exactly $j_1$ occurrences of $h_{u_1}$ and $j_2$
occurrences of $h_{u_2}$, for crossing positions $u_1,u_2$ in the two windows,
and with the remaining window terms weighted by the complex variables $Z$.
Between consecutive occurrences of the markers $h_{u_1}$ and $h_{u_2}$ such a
term therefore contains powers of the weighted ordinary sum
\begin{align}
 G_Z=z_MH_M+\sum_{u}z_uh_u ,
 \label{eq:alvv16_cmb_ordinary}
\end{align}
where $u$ runs over the window positions other than $u_1,u_2$, together with
the terms of $H_L$ and $H_R$.

Fix a cut inside $J_M$. Every window term lies entirely on one side of it, so
the left windows contribute one operator of the form $X\otimes 1$ and the right
windows one of the form $1\otimes Y$. With $H_M$ of operator Schmidt rank at
most $\tilde B$ at the cut, the factor (\ref{eq:alvv16_cmb_ordinary}) has
operator Schmidt rank at most $\tilde B+2$ there, the markers have rank one,
and submultiplicativity bounds a product of at most $k$ such factors by
\begin{align}
 (\tilde B+2)^k .
 \label{eq:alvv16_cmb_bound}
\end{align}
Expanding the operators on the qudits of the two windows inside
$\{i_1,\ldots,i_2\}$ in a product basis, the step of the proof that multiplies
the number of operators by at most $d^{4\ell}$, contracts the window qudits
with fixed configurations and does not increase the rank at a cut inside
$J_M$. The printed argument, read in this way, gives the bound
(\ref{eq:alvv16_cmb_bound}) at every cut inside $J_M$. It gives no argument
that removes the two summands contributed by the windows, and whether the
printed bound $\tilde B^k$ holds is open. The bound
(\ref{eq:alvv16_cmb_bound}) is not a refutation of the proposition, and it is
at most $(2\tilde B)^k$ for $\tilde B\geq2$.

\section{What is formalized}

The formalization reads an MPO for $H_M$ as a site-dependent trace contraction
of tensors of virtual dimension $B$ with a closing matrix. Its common middle
family for every polynomial of degree at most $k$,
\leanid{TNLean.MPS.AreaLaw.exists_common_physical_polynomial_mpo_family_fin},
has the source's count shape, with at most $D^2$ members for
$D=(qk)^{18\ell+2\lfloor k/\ell\rfloor}$, and the members act on the fixed
middle interval. Each member has open-boundary bond dimension at most
\begin{align}
 \max(1,B)^{2k},
 \label{eq:alvv16_cmb_formal_bond}
\end{align}
and a trace presentation of virtual dimension at most $\max(1,B)^k$. The
square in (\ref{eq:alvv16_cmb_formal_bond}) comes from the bottom of the
chain. There, the selected-coefficient family
\leanid{TNLean.MPS.AreaLaw.exists_common_physical_twoMarkerCoefficient_mpo_family}
splits each middle word by its number $t$ of factors of $H_M$ into members
that are trace contractions of the $t$-fold stacked tensor with a closing
matrix of size $B^t$. Reading such a member in open-boundary form costs the
square $(B^t)^2$, because the closing matrix collects the summed
configurations of the window qudits and has no small rank in general. The
square is carried unchanged through the selected-middle family, the
fixed-middle family, the crossing-frequency grid and the polynomial family,
and the efficient soft families state their member bond in the form
(\ref{eq:alvv16_cmb_formal_bond}).

\section{Plan}

The bound (\ref{eq:alvv16_cmb_bound}) can be reached by re-plumbing the
existing chain, without a parallel one.
\begin{enumerate}
 \item Replace the members at the bottom by the blocks
  $\Pi_\eta\Psi_w\Pi_\theta$ of each middle word $\Psi_w$, where $\eta,\theta$
  run over the configurations of the window qudits inside the middle interval
  and $\Pi_\eta$ is the corresponding product of one-site projections. Such a
  block has operator Schmidt rank one at a cut among the window qudits and at
  most the rank of $\Psi_w$ at a cut inside the strict core. The split by the
  number of factors of $H_M$ is no longer needed, so the factor $k+1$ leaves
  the count.
 \item Strengthen the prescribed-middle decomposition so that the ordinary
  sum on the window qudits is given as $K_L+K_R$, with $K_L$ supported on the
  left windows and $K_R$ on the right ones. At present only its support in
  the union of the two is recorded.
 \item Take the input as an open-boundary MPO, or directly as an operator
  Schmidt rank bound $\tilde B$ for $H_M$ at every cut, and bound the rank of
  $\Psi_w$ at a cut inside the strict core by (\ref{eq:alvv16_cmb_bound}) using
  the existing submultiplicativity of operator Schmidt rank.
 \item Restate the bond clause of the four upper levels with
  $(\tilde B+2)^k$ in place of (\ref{eq:alvv16_cmb_formal_bond}), and adjust
  the efficient soft families that consume it.
\end{enumerate}
The change touches about six modules and is estimated at 1500--2500 lines.
With it, the proposition would be formalized with the bond
(\ref{eq:alvv16_cmb_bound}) as a local fix of the printed $\tilde B^k$.

\section{Verdict}

Scope restriction, open. The existence of a common middle family with the
source's count is formalized with the bond (\ref{eq:alvv16_cmb_formal_bond}).
The printed proof supports the bond (\ref{eq:alvv16_cmb_bound}) at cuts inside
$J_M$; whether $\tilde B^k$ holds is open. The re-plumbing of the plan is not
started.

\bibliographystyle{plain}
\bibliography{references}

\end{document}
